%==============================================================================
%  Appendix G --- which split each operator picks.
%  Filename carries the topic, not the letter (see the master's header note).
%  Every number here comes from analysis/supporting/eta_by_operator.ipynb.
%==============================================================================

\section{Which Split Each Operator Picks}\label{app:eta}

Two cohorts of 1000 trees, Kingman and birth--death, at $m=1000$ and $L=10^4$ under
\Cref{asm:gtr,asm:clock}, with no sub-sampling. Each operator uses the rounding rule its
own theorem states: $k$-means on the Fiedler vector of $L(S)$, sign on the
leading-magnitude eigenvector of $\Bop=H\Dmat H$. Per tree we record the imbalance of the
returned split, $\eta=\max(n_1,n_2)/\min(n_1,n_2)$, and whether that split is an edge of
the true tree.

\begin{figure}[htbp]
    \centering
    \begin{subfigure}[t]{0.49\textwidth}
        \captionsetup{singlelinecheck=false, justification=raggedright, skip=1pt}
        \subcaption{Kingman, $L(S)$, $k$-means}\label{fig:eta-king-ls}
        \includegraphics[width=\textwidth]{eta_hist_kingman_S.png}
    \end{subfigure}
    \hfill
    \begin{subfigure}[t]{0.49\textwidth}
        \captionsetup{singlelinecheck=false, justification=raggedright, skip=1pt}
        \subcaption{Kingman, $\Bop$, sign}\label{fig:eta-king-b}
        \includegraphics[width=\textwidth]{eta_hist_kingman_B.png}
    \end{subfigure}

    \vspace{3mm}

    \begin{subfigure}[t]{0.49\textwidth}
        \captionsetup{singlelinecheck=false, justification=raggedright, skip=1pt}
        \subcaption{birth--death, $L(S)$, $k$-means}\label{fig:eta-bd-ls}
        \includegraphics[width=\textwidth]{eta_hist_bd_S.png}
    \end{subfigure}
    \hfill
    \begin{subfigure}[t]{0.49\textwidth}
        \captionsetup{singlelinecheck=false, justification=raggedright, skip=1pt}
        \subcaption{birth--death, $\Bop$, sign}\label{fig:eta-bd-b}
        \includegraphics[width=\textwidth]{eta_hist_bd_B.png}
    \end{subfigure}
    \caption{\textbf{Imbalance of the returned split, 1000 trees per model.} Green is a
    genuine single-edge bipartition of the true tree, red is not; dashed line is the
    median. All four panels share bins and axes.}
    \label{fig:eta-hist}
\end{figure}

\begin{table}[htbp]
\centering
\caption{Both operators on both cohorts, 1000 trees each.}
\label{tab:eta-agreement}
\begin{tabular}{lrrrr}
\toprule
& \multicolumn{2}{c}{Kingman} & \multicolumn{2}{c}{birth--death}\\
\cmidrule(lr){2-3}\cmidrule(lr){4-5}
& $L(S)$ & $\Bop$ & $L(S)$ & $\Bop$\\
\midrule
valid single-edge split & $996$ & $986$ & $996$ & $\mathbf{555}$\\
median $\eta$           & $3.00$ & $1.85$ & $3.08$ & $1.53$\\
maximum $\eta$          & $499$  & $24.6$ & $999$  & $5.3$\\
\midrule
only $L(S)$ valid       & \multicolumn{2}{r}{$14$}  & \multicolumn{2}{r}{$443$}\\
only $\Bop$ valid       & \multicolumn{2}{r}{$4$}   & \multicolumn{2}{r}{$2$}\\
identical $\eta$        & \multicolumn{2}{r}{$64.1\%$} & \multicolumn{2}{r}{$33.2\%$}\\
\bottomrule
\end{tabular}
\end{table}

\noindent Four things follow.

\begin{enumerate}[leftmargin=2em, itemsep=3pt, topsep=3pt]
\item \textbf{On Kingman trees the operators differ in \emph{which} edge, not in whether
      they find one} --- $99.6\%$ and $98.6\%$ valid, and the same split on $64.1\%$ of
      trees.
\item \textbf{$\Bop$'s output is confined to near-balanced splits}: $\eta\le24.6$ on
      Kingman and $\eta\le5.3$ on birth--death, against true splits reaching $\eta=999$.
      This is \Cref{lem:dist-edge} --- each edge contributes at least
      $2\tau_e n_1(e)n_2(e)/m$ to $|\lambda_1(\Bop)|$, and $n_1(e)n_2(e)$ is maximal at
      balance. Where the rules disagree, $\Bop$'s split is the more balanced one in $355$
      of $359$ Kingman cases and $654$ of $668$ birth--death cases.
\item \textbf{That confinement breaks $\Bop$ on imbalanced trees.} Its validity falls from
      $98.6\%$ to $55.5\%$ while $L(S)$ holds at $99.6\%$; on $443$ birth--death trees
      $L(S)$ returns an edge and $\Bop$ does not. Lacking a near-balanced edge it does not
      decline --- it returns a near-balanced \emph{non}-edge.
\item \textbf{A narrow $\eta$ distribution is a property of the operator, not of the
      tree.} Comparing the two routes' sampling rates is meaningful only with the target
      split held fixed; otherwise the rates describe different recovery problems, and on
      birth--death trees one is conditioned on a split absent from the tree.
\end{enumerate}
